\scanpage{67}
\begingroup\makeatletter\tagsleft@true\makeatother
\uline{2°} On obtient $Z^2=\dfrac{1}{k^2}\bigl[\unsure{p|\wp}(z-z_0)-e_1\bigr]$
\[
\left\{\begin{aligned}
e_1&=-\frac{1+k^2}{\deleted{\unsure{3}}3}\\
e_3&=\frac{2k^2-1}{3}\\
e_2&=\frac{2-k^2}{3}
\end{aligned}\right. \tag*{(2)}
\]

La solution envisagée est
\[
\boxed{\operatorname{sn}^2 z=\frac{1}{k^2}\bigl[p(z-\omega_1)-e_1\bigr]} \tag*{(3)}
\]

\uline{3°} Avec la même fonction elliptique, on a une solution générale $Z^2=\dfrac{1}{p(z-z_0)-e_1}$ et la solution particulière est
\[
\boxed{\operatorname{sn}^2 z=\frac{1}{pz-e_1}} \tag*{(4)}
\]

On montrerait facilement l'uniformité de \deleted{$Z(z)$} $\operatorname{sn}z$ (les racines de $pz-e_1=0$ étant doubles)

$\ast\ast\ast$ Ces deux formules impliquent une identité :
\[
(pz-e_1)(p(z-\omega_1)-e_1)=k^2
\]
facile à vérifier, en passant à la fonction $\sigma$, ou en montrant que le premier mb qui est elliptique, n'a pas de pole, et est donc constant -- ceci est vrai pour \deleted{tte} fonction $pz$ -- On trouve d'ailleurs incidemment la formule $\dfrac{e^{2\eta_1\omega_1}}{\sigma^{\unsure{4|2}}\omega_1}=3e_1^2-\dfrac{g_2}{4}$, \deleted{qui permet le calcul de $\eta_1$ et \unsure{à|de} fort. $\eta_2$ et $\eta_3$}

Dans le cas actuel, on a $k^2=\dfrac{e^{2\eta_1\omega_1}}{\sigma^4\omega_1}=3e_1^2-\dfrac{g_2}{4}$ $\,-\ast\ast\ast$

Toujours est il que nous avons $\operatorname{sn}^2(z+\omega_1)\operatorname{sn}^2z=\dfrac{\deleted{\unsure{1}}1}{k^2}$ ; si nous choisissons $\omega_1$ et $k$ convenablement nous aurons
\[
\boxed{\begin{aligned}
&\operatorname{sn}(z+\omega_1)\operatorname{sn}z=\frac{1}{k}\\[-.5ex]
&\scriptstyle\unsure{=}
\end{aligned}} \tag*{(5)}
\]
\endgroup
% archive sheet number: 66
